\section{Related Work}
\label{sec:related_work}

Action recognition literature has progressed across three distinctive architectural eras: handcrafted spatiotemporal descriptors, deep convolutional neural networks, and attention-based video vision transformers.

\subsection{Classical and Convolutional Video Architectures}
Early action recognition techniques relied on handcrafted space-time interest points (STIP) and optical flow trajectories~\cite{rodriguez2008action}. With the rise of deep learning, two-stream convolutional networks decoupled video processing into an appearance stream operating on RGB frames and a temporal stream operating on pre-computed optical flow fields. To bypass the prohibitive computational cost of dense optical flow extraction, two-stage CNN-RNN architectures were introduced, wherein 2D CNNs (e.g., ResNet, VGG) extracted frame-level feature vectors that were sequentially processed by LSTMs or GRUs~\cite{donahue2015long}. While efficient, these recurrent models fail to model cross-frame spatial deformations jointly.

To capture unified spatiotemporal structures, 3D Convolutional Networks (3D CNNs) were developed. Carreira and Zisserman introduced Inflated 3D ConvNets (I3D)~\cite{carreira2017i3d}, which inflated 2D ImageNet filters into 3D kernels. Feichtenhofer \etal{} introduced SlowFast Networks~\cite{feichtenhofer2019slowfast}, which employ two parallel pathways operating at differing frame rates: a Slow pathway with high spatial capacity capturing fine spatial semantics, and a Fast pathway with high temporal refresh rate capturing motion dynamics. To address the computational burden of 3D CNNs, Feichtenhofer proposed X3D (Expanding Architectures)~\cite{feichtenhofer2020x3d}, which systematically expands a tiny 2D image baseline across spatial, temporal, and channel dimensions to maximize accuracy per FLOP. Despite their efficiency, 3D CNNs remain constrained by local convolutional receptive fields, limiting their capacity to model long-range temporal reasoning.

\subsection{Video Vision Transformers}
The breakthrough of the Vision Transformer (ViT)~\cite{dosovitskiy2020vit} in 2D image recognition inspired the development of video transformers. Arnab \etal{} introduced ViViT~\cite{arnab2021vivit}, exploring factorized spatial and temporal self-attention across video token sequences. Concurrently, Bertasius \etal{} presented TimeSformer~\cite{bertasius2021timesformer}, establishing that divided space-time attention achieves superior computational scaling compared to joint spatiotemporal attention. Liu \etal{} introduced the Video Swin Transformer~\cite{liu2022videoswin}, which generalizes shifted window self-attention to 3D video volumes, creating hierarchical multiscale representations suitable for dense video tasks. However, training video transformers from scratch or fine-tuning supervised backbones on moderate-sized datasets frequently induces severe overfitting due to the absence of inductive convolutional bias.

\subsection{Self-Supervised Masked Video Modeling}
To alleviate the annotation bottleneck and prevent visual shortcut learning, masked autoencoding was extended to the temporal domain. Building upon the Masked Autoencoder (MAE) paradigm of He \etal{}~\cite{he2022masked}, Tong \etal{} developed VideoMAE~\cite{tong2022videomae}. VideoMAE incorporates tubelet patch embedding ($16 \times 16 \times 2$) and introduces an aggressive masking strategy (up to 90\% of tokens masked out). Due to the high temporal redundancy between consecutive video frames, high-ratio masking prevents trivial pixel interpolation from adjacent frames, compelling the autoencoder to learn high-level motion kinematics and spatiotemporal representations. Wang \etal{} further scaled this paradigm in VideoMAE V2~\cite{wang2023videomaev2} using dual masking. In this work, we harness VideoMAE ViT-B as the foundational representation engine for fine-grained sports action classification.
